% Metódy inžinierskej práce
\documentclass[10pt,twoside,english,a4paper]{article}

\usepackage[english]{babel}
\usepackage[T1]{fontenc}
\usepackage[utf8]{inputenc}
\ifdefined\XeTeXversion
\usepackage{fontspec}
\fi
\usepackage{url}
\usepackage[hidelinks]{hyperref}
\usepackage{cite}

\pagestyle{headings}


\newcommand{\projtitle}{Knowledge, Inclusion \& Real-Time Assistance}
\newcommand{\projacronym}{KIRTA}

\title{\projacronym: \projtitle\thanks{Semester project in Metódy inžinierskej práce, academic year 2026/27, supervisor: Ing. Vladimír Mlynarovič, PhD.}}

\author{Maksym Bulat, Yevhenii Demianchuk, Yehor Donin\\[2pt]
	{\small Slovak University of Technology in Bratislava}\\
	{\small Faculty of Informatics and Information Technologies}\\
	{\small \texttt{xbulat@stuba.sk, xdemianchuk@stuba.sk, xdonin@stuba.sk}}
	}

\date{\small \today}

% added this, cause defoult name is 'References' /this command rename the bibliography heading into Literature
\addto\captionsenglish{\renewcommand{\refname}{Literature}}


\begin{document}

\maketitle

\section*{Annotation}

\subsection*{The essence of the project}

\projacronym{} is a voice-first daily-life AI companion for blind and low-vision people.

\subsection*{Target audience}

More than 1 billion people experience sight loss; 43 million of them are blind~\cite{iapb}. \projacronym's primary target audience is blind and low-vision people. However, later expansion into people with other disabilities is possible.

\subsection*{Problem}

Nowadays, people depend on smartphones for navigation, communication, shopping, studying, working, banking, entertainment, and other mundane tasks. People with blindness are no exception; for them, smartphones serve as a lifeline~\cite{senjam2022}, replacing their sight. Yet using a phone requires them to hold it, orient themselves on the touchscreen, and ask other people for help when needed. Existing apps mostly work on request; for example, reading text from a photo/screen one at a time or helping with navigation. They also struggle with understanding users' intent~\cite{xie2025}.

\subsection*{Solution}

Instead of holding the phone all the time, the phone is secured by a case with a strap. Its camera acts as the user's eyes, analyzing the surroundings with AI. Then the information is narrated to the user's headphones. Most of the time, the app is silent, breaking the silence only in critical moments, when the user is in danger, or on request. \projacronym{} is also expected to help with navigation, calls, messages, and alerts. More functions are to come later; for example, ordering goods and food, consent-based face recognition~\cite{gdpr, aiact, akter2022}, tourist guidance, and integration into desktop operating systems. Security comes first, so everything involving biometric data and hazard detection will run locally, and payments are always confirmed by the user in the service's own app. Other features use the cloud when the internet is available.

\subsection*{Monetization}


We propose free local assistance and an optional paid cloud-description plan, with pilot funding and sponsored seats to be explored; WHO identifies affordability as a barrier to assistive technology~\cite{whotechnology2024}. The budget and revenue assumptions are detailed below.

\section*{List of functions}
\subsection*{Video-to-audio translation}
\subsection*{Navigation}

The navigation design for KIRTA combines GPS with local camera analysis. GPS provides the general position on the map and supports the route to a destination. The local AI follows the camera feed and identifies what is happening near the person. These sections describe a preliminary design, with implementation and testing to follow.

A person may be on the correct route and still face a blocked pavement, stairs or a vehicle approaching from the side. The navigation feature needs to deal with that part of the journey too. The phone will be attached to a strap, with its camera facing the surroundings, leaving the person's hands free.

GPS.gov reports typical smartphone accuracy within about 4.9 metres under open sky. Accuracy gets worse near buildings, bridges and trees~\cite{gpsaccuracy}. GPS gives the app a starting position, while the camera supplies a view of the immediate surroundings. Another goal is to use visual information to refine the position estimate as the person moves. Matching what the camera sees to a particular place on the map remains an open technical question.

Hazard detection is planned to run locally on the phone, analysing the camera feed throughout the walk. The recognition tasks include stairs, obstacles along the path and moving vehicles. With vehicles, movement matters as much as recognising the object: the system needs to assess whether it is approaching the user's path. Evaluation will measure detection errors and the delay between a hazard appearing and a warning being spoken.

In the study of Be My AI by Xie and colleagues, participants described stopping to take pictures and checking whether they had captured the information needed for navigation. Some descriptions also lacked details about nearby obstacles~\cite{xie2025}. These findings are relevant to KIRTA's use of a camera feed that updates as the user moves.

NavCog offers another useful reference. It used Bluetooth beacons installed in the environment and was tested with six blind participants on a university campus~\cite{navcog2016}. Our design centres on the camera carried by the user. Soundscape is relevant to the audio side: it explored spatial sound for nearby places and a chosen destination~\cite{soundscape}. These projects give us examples of both location support and ways to communicate directions.

For everyday use, KIRTA should stay quiet until there is useful information to pass on. A possible hazard takes priority; a fuller description of the surroundings can be requested. The assistance is intended to work alongside the person's mobility skills and aids, including a white cane or guide dog.

\subsection*{Narration of on-screen activity}

When a new SMS arrives, KIRTA should let the user know and offer to read it. This is the simplest example of screen narration in the project. The same feature covers incoming calls, notifications and changes in the open app. AI is intended to interpret this information and produce a short explanation that can be spoken through the headphones.

The explanation needs to fit what the person is doing. After pressing a button, the useful information is usually the result: whether the action succeeded, whether the app is still waiting or whether an error appeared. On opening an unfamiliar screen, a brief description of its contents and available actions may be more useful. These situations call for different amounts of detail. Repeating the whole screen after every change would make even a simple task tiring.

Android provides two possible sources for this feature. Accessibility services can inspect available interface elements and receive events such as clicks and changes of focus~\cite{androidaccessibility}. A notification listener service receives updates when notifications are posted or removed~\cite{androidnotifications}. Access depends on the required configuration and permissions. In the initial Android design, these mechanisms would supply information for the AI to explain. Image-based recognition is a further option for screens where the exposed text and controls are insufficient.

W3C guidance helps specify the information worth keeping in that explanation. An accessible control has a name and role, together with relevant state information. A checkbox needs its label and whether it is selected~\cite{wcagname}. Status messages describe progress, errors and action results without requiring the user to move focus to the message~\cite{wcagstatus}. A meaningful focus order also helps people follow web content~\cite{wcagfocus}. We use these web accessibility principles as a reference for the spoken interaction in KIRTA.

Our proposed interaction gives the user control over longer descriptions and private messages. A short notification can be followed by a request to read more, repeat it or stop speaking. The account of AI errors in Xie and colleagues' study is relevant here too: generated explanations can include details that the input does not support~\cite{xie2025}. Screen narration therefore needs to stay grounded in the available interface information. Confirmation that a message was sent, for example, must come from the app's reported state.

Phone narration will share the audio output with navigation warnings. The proposed design gives a possible hazard outside priority over an ordinary notification and avoids repeated announcements of unchanged information. Evaluation should include tasks such as noticing a new message and understanding a change on the screen. The records should include missed updates, incorrect descriptions and unnecessary interruptions.

\subsection*{Friend}

\subsection*{Voice command recognition}

Voice commands will be one of the main ways to use KIRTA. For a blind person, it is often easier to say what they need than to look for a button on a touchscreen. The user could say ``Describe what is in front of me'', ``Start navigation'' or ``Call Anna''. Our goal is to make these commands short and easy to remember. The application should not require the user to say a long, exact sentence every time.

To recognize speech, we could use the Android \texttt{SpeechRecognizer} API. It receives spoken words and returns text which can then be checked by the application~\cite{androidspeech}. At first, KIRTA could support about ten basic commands. For example, if the recognized text contains ``call'' and a contact name, the application would understand that the user wants to make a phone call. Other commands could request a description of the camera image, repeat the last answer or stop the audio output.

Recognition will not always work perfectly. Traffic, other people talking and different accents can make words harder to understand. If KIRTA is not sure what the user said, it should ask for the command again. Some Android phones support on-device speech recognition, but it is not available for every phone or language~\cite{androidspeech}. We would prefer local processing when possible because it can also help with privacy. For testing, we could try the same ten commands in a quiet room and outside, and count how many times they are understood correctly. We could also ask a few classmates to try the same commands. It would help us see whether the words are clear enough for people who did not write the program.

\subsection*{Voice-controlled calls and messages}

Another useful function is starting calls and writing messages without searching through the phone menus. For example, a user could say ``Call Mom''. KIRTA would search for the contact and repeat the name so the user knows who was selected. The same idea works for messages: ``Send a message to Peter'' would start a short conversation where KIRTA asks what the message should say. The user could then dictate the text normally.

Before doing anything important, KIRTA should confirm what it understood. It could say ``Message to Peter: I will be late. Is that correct?'' and wait for ``yes'' or ``no''. This is especially useful when two contacts have similar names or when speech recognition makes a mistake. The user should also be able to cancel the action at any point.

For a first version, we do not need to build our own phone or messaging service. Android provides \texttt{Intent} actions for opening the phone dialer and preparing an SMS message~\cite{androidintents}. The normal Android application would finish the call or send the message. This also means that some final steps may still depend on the accessibility functions of the phone, such as TalkBack. Later, we could look into other messaging applications, but starting with calls and SMS is a more realistic plan.

One problem could be finding the correct contact. People often save more than one number under the same name, for example a personal and a work number. In this case, KIRTA should ask which number the user wants. We would also try messages with short names, longer names and names that sound similar. The aim is not to make the phone call faster at any cost, but to make the action understandable.

\subsection*{Programming and technologies}

We plan to develop KIRTA as an Android application. We would use Kotlin because it is supported for Android development and can be used in Android Studio~\cite{androidkotlin}. We chose Android mainly because smartphones already have a camera, microphone, speaker and support for Bluetooth headphones. We can use these parts instead of creating separate hardware for the first prototype.

The application will need a few basic parts that work together. The first part listens to the user and recognizes commands. The second part decides what the command means and which action to start. A third part reads answers aloud using Android \texttt{TextToSpeech}~\cite{androidtts}. For example, after hearing ``What is in front of me?'', the program would request information from the camera part and then speak the result. The sound could be played through connected Bluetooth headphones using the phone's normal audio settings.

For camera access, we could use the CameraX library. It is made for Android camera applications and could later provide images to an existing recognition model~\cite{androidcamerax}. We do not plan to train a new AI model from zero for this stage. Our first goal would be to connect simple functions and make them work reliably. The other team members could then work on navigation and recognizing possible dangers, while the voice interface would connect these features to the user.

We would test the first version with a small set of commands, different contact names and simple spoken messages. We also need to check permissions for the microphone, camera and contacts. It is possible that the app will work differently on different Android devices, so testing on a real phone will be important. At this stage, KIRTA is still a project idea, and the exact technologies may change after the first experiments.

A simple development plan would start with one Android screen and a button for recording a command. Next, we could add spoken answers and a few actions, such as opening the dialer. Only after these functions work would we connect the camera and other features from the team. It is easier to find problems step by step than to build everything at once.





\clearpage
\section*{Budget and monetization}

\textbf{Scope and assumptions.} We propose a three-month Android pilot for 100 users; A denotes our allowances, with citations identifying benchmarks or activities. We value 300 contributed development hours using Eurostat's 2025 EU-wide employer-cost benchmark, EUR~34.9/hour, rather than a software-developer quotation~\cite{eurolabour2025}. Hardware and testing follow the smartphone architecture and WHO's user-involvement guidance~\cite{androidcamerax,whotechnology2024}.

\begin{center}
\small
\renewcommand{\arraystretch}{1.12}
\begin{tabular}{@{}p{3.3cm}p{6.3cm}r@{}}
\hline
\raggedright \textbf{Pilot item} & \raggedright \textbf{Calculation and evidence} &  \textbf{EUR} \tabularnewline
\hline
\raggedright Development (A) & \raggedright $3\times100$ hours $\times34.9$~\cite{eurolabour2025} &  10,470.00 \tabularnewline
\raggedright Hardware (A) & \raggedright $2\times200$ for phones; 100 for straps/headsets~\cite{androidcamerax} &  500.00 \tabularnewline
\raggedright Accessible testing (A) & \raggedright $10\times25$ participant allowances; $10\times34.9$ facilitation~\cite{whotechnology2024,eurolabour2025} &  599.00 \tabularnewline
\raggedright Privacy/qualification review (A) & \raggedright $20\times34.9$; preliminary allowance requiring a professional quote~\cite{gdpr,mdcg2019software,eurolabour2025} &  698.00 \tabularnewline
\raggedright Cloud inference (A) & \raggedright 100 users, 90 days, 30 requests/day; calculation below~\cite{geminipricing2026,ecbrate20261030} &  65.25 \tabularnewline
\raggedright Hosting/logs (A) & \raggedright $3\times20$ monthly allowance; actual usage billed separately~\cite{cloudrunpricing2026} &  60.00 \tabularnewline
\raggedright Play registration & \raggedright One-off USD~25, converted at the ECB planning rate~\cite{playregistration2026,ecbrate20261030} &  22.38 \tabularnewline
\hline
\raggedright Resource subtotal & \raggedright Authors' sum of the above items &  12,414.63 \tabularnewline
\raggedright Contingency (A) & \raggedright 20\% of the subtotal &  2,482.93 \tabularnewline
\raggedright \textbf{Total resources} & \raggedright \textbf{Including contributed development time} &  \textbf{14,897.56} \tabularnewline
\hline
\end{tabular}
\end{center}

\textbf{Cloud calculation.} For a candidate model, Google's published Gemini~2.5 Flash-Lite standard paid rates are USD~0.10/million input tokens and USD~0.40/million output tokens; assuming 1,500 input tokens including an image and 300 total billable output tokens per request gives $270{,}000(1500\times0.10+300\times0.40)/10^6=\mathrm{USD}~72.90$ for the pilot~\cite{geminipricing2026}. At the ECB's 30 October 2026 reference rate, EUR~1 = USD~1.1172, this is EUR~65.25 before supplier taxes and exchange fees; real image sizes, retries and model changes require metering and per-user limits~\cite{ecbrate20261030,geminipricing2026}. With development contributed in kind, the base cash request including contingency is approximately EUR~2,335; commercial certification, insurance and a production launch require a separate budget following software qualification~\cite{eurolabour2025,mdcg2019software}.

\textbf{Revenue proposal and viability.} We propose free local warnings and basic voice/accessibility functions, with an optional EUR~4.99/month cloud-description plan capped at 900 requests; this responds to the affordability barrier identified by WHO~\cite{whotechnology2024}. For an illustrative Slovak sale, assume standard 23\% VAT and Google Play's EEA subscription charges of 10\% service plus 5\% billing: $4.99/1.23\times0.85-0.30-0.80\approx\mathrm{EUR}~2.35$ contribution per subscriber-month, where EUR~0.30 inference and EUR~0.80 support are our allowances~\cite{slovakvat2026,playfees2026,geminipricing2026}. Recovering EUR~14,900 over twelve months plus EUR~160/month for hosting and four maintenance hours would require about 600 paying subscribers; this is a sensitivity scenario, not evidence of demand~\cite{eurolabour2025,cloudrunpricing2026}. Sponsored seats and pilot grants remain uncommitted; the European Accessibility Act establishes accessibility requirements rather than funding entitlement~\cite{eaa}.

\clearpage
\section*{Legal and medical challenges}

\textbf{Personal data and bystanders.} Under GDPR Articles~5, 6, 25 and 32, we would document lawful bases and minimise camera, location, message and contact data through local processing, short-lived buffers, minimal permissions and encryption~\cite{gdpr}. Before cloud use, we would specify roles, retention, deletion, Article~28 processor contracts and Chapter~V transfer safeguards, and assess Article~35 DPIA necessity~\cite{gdpr}. Unique face identification attracts Article~9 protection, and wearer consent cannot cover bystanders; we would disable face recognition pending a lawful design and assessment against EDPB video guidance~\cite{gdpr,edpbvideo2020}.

\textbf{Medical-device and AI classification.} The MDR includes software intended to compensate for disability; avoiding diagnosis or adding a disclaimer does not establish exemption, so we would document each module's intended purpose and seek qualification/classification advice using MDCG~2019-11~\cite{mdr2017,mdcg2019software}. AI Act Article~6(1) requires both an Annex~I-covered product/safety component and third-party conformity assessment; we would assess these conditions and the separate Annex~III route before classifying KIRTA~\cite{aiactarticle62026}. Applicable medical-device risk management, technical documentation and clinical evaluation would require funding beyond this pilot~\cite{mdr2017,mdcgclinical2020}.

\textbf{Store access and user control.} Google Play reserves accessibility-tool status for apps primarily supporting people with disabilities and requires review; our AccessibilityService declaration, store description and demonstration would establish this purpose~\cite{playaccessibility2026}. We would offer spoken permission explanations, cancellation and confirmation before calls/messages; Android permissions alone do not establish GDPR compliance~\cite{gdpr,androidintents}.

\textbf{Navigation injury and technical failure.} Missed obstacles, delayed warnings and an obstructed camera are foreseeable hazards; applying MDR Annex~I risk-management and software-reliability principles as a design benchmark, we propose a hazard log, offline warning priorities and explicit failure announcements for camera, battery or processing loss~\cite{mdr2017}. KIRTA would supplement established mobility aids, and would never certify that a road is safe to cross; the aim is additional information within an individually suitable combination of assistive products~\cite{whotechnology2024,mdr2017}.

\textbf{Hearing, attention and AI errors.} WHO identifies loudness and exposure duration as hearing-risk factors; our proposed controls are adjustable low-volume speech, short warnings and user-selected audio arrangements, with environmental audibility checked in supervised trials~\cite{wholistening2026}. WHO also warns that multimodal models can generate false statements and encourage automation bias; KIRTA should announce uncertainty, avoid unsupported assertions and exclude diagnosis, treatment and medication decisions from the pilot~\cite{who2024ai}. These restrictions guide design and do not demonstrate clinical safety~\cite{mdcgclinical2020}.

\textbf{Evidence before deployment.} Evaluation would progress from recorded scenes to indoor routes and supervised outdoor trials, using accessible consent and stopping rules, and recording missed hazards per obstacle, false warnings/hour, 95th-percentile latency, task completion and workload across lighting, noise and camera positions~\cite{mdcgclinical2020}. Users and mobility specialists would agree acceptance thresholds; MDCG~2020-1 separates technical and clinical performance, so a small usability pilot would not establish clinical benefit or readiness for unsupervised traffic use~\cite{mdcgclinical2020}.

\clearpage
{\fontsize{8.5}{10}\selectfont
\bibliography{literatura}
}
\bibliographystyle{unsrt}

\end{document}
